\chapter{ANÁLISE DOS RESULTADOS}
\label{capitulo4}

Apresente neste capítulo os resultados obtidos na pesquisa de maneira objetiva, clara e lógica. Os dados apresentados devem ser interpretados e confrontados com a literatura científica revisada no Capítulo~\ref{capitulo2}.

\section{APRESENTAÇÃO DOS DADOS E RESULTADOS}
\label{sec:apresentacao_resultados}

Apresente os resultados preliminares ou agregados decorrentes da aplicação do método descrito anteriormente. Utilize gráficos, quadros e tabelas para organizar as informações de maneira didática e concisa.

\section{ANÁLISE ESTATÍSTICA E DISCUSSÃO DOS ACHADOS}
\label{sec:analise_discussao}

Discuta o significado dos resultados à luz das hipóteses formuladas no Capítulo~\ref{capitulo1}:
\begin{itemize}
  \item Avalie se as hipóteses iniciais foram confirmadas ou refutadas pelos dados empíricos;
  \item Destaque relações estatisticamente significativas ou comportamentos anômalos encontrados;
  \item Compare os resultados alcançados com os estudos prévios de referência citados na literatura.
\end{itemize}
